% ECONOMICS_PROBLEM: {"id": "F34-105", "title": "Restructuring with diverse creditors", "jel_code": "F34", "source_id": "OP-0256"}
% !TeX program = xelatex
\documentclass[11pt,a4paper]{article}
\usepackage[margin=23mm]{geometry}
\usepackage{fontspec}
\setmainfont{Latin Modern Roman}
\usepackage{amsmath,amssymb,mathtools}
\usepackage{xcolor,enumitem,booktabs,longtable,array,xurl,hyperref}
\definecolor{muted}{HTML}{53636C}
\hypersetup{colorlinks=true,urlcolor=blue,linkcolor=blue}
\setlength{\parindent}{0pt}
\setlength{\parskip}{5pt}
\newcommand{\R}{\mathbb R}
\newcommand{\E}{\mathbb E}
\newcommand{\N}{\mathbb N}
\newcommand{\ind}{\mathbf 1}
\newcommand{\Var}{\operatorname{Var}}
\newcommand{\Cov}{\operatorname{Cov}}
\newcommand{\supp}{\operatorname{supp}}
\DeclareMathOperator*{\argmin}{arg\,min}
\DeclareMathOperator*{\argmax}{arg\,max}
\newcommand{\entrymeta}[3]{{\sffamily\small\color{muted}\textbf{\texttt{#1}}\quad|\quad #2\par #3\par}}
\newcommand{\Problemref}[1]{\texttt{#1}}
\newcommand{\JELref}[2]{\texttt{#2} (JEL \texttt{#1})}
\begin{document}
\section*{Restructuring with diverse creditors}
\textbf{Problem ID:} \texttt{F34-105}\quad \textbf{JEL:} \texttt{F34}\par
% BEGIN ECONOMICS STATEMENT
\label{problem:OP-0256}
{\sffamily\small\color{muted}Original subject group: Fiscal policy and sovereign debt\par}
\label{definitions:problem:30007537}
\textbf{Audit classification:} Published research agenda. Official-creditor implications are an agenda; the multilateral mechanism is editorial and overlaps game theory. Verification concerns the cited version, not first priority or proof that this exact editorial specification remains unsolved on October 8, 2026.
\subsection*{Definitions and precise research target}
\paragraph{Editorial mathematical specification.} Fix an insolvent sovereign owing claims \(b_j>0\) to a finite set of creditors \(j\), including domestic banks, private foreign investors, multilateral institutions, and official bilateral lenders. Each creditor has a declared seniority, outside option, legal enforcement capacity, and participation constraint. The sovereign has stochastic repayment capacity \(R_t\) and a real social cost from delayed settlement. A restructuring mechanism chooses repayments \(x_{jt}\), creditor voting rules, and access to new financing after reports about capacity and creditor types.
Find mechanisms minimizing expected discounted crisis costs subject to aggregate feasibility
\[
\sum_j x_{jt}\le R_t+\text{permitted new financing at }t,
\]
creditor participation, any statutory seniority constraints, and incentive compatibility for specified private information. Characterize when participation by diverse official and private creditors permits timely settlement, and when holdouts or incompatible legal claims make the target allocation unattainable. The mechanism must include how anticipated restructuring terms affect initial lending and borrowing; a one-period negotiated haircut is insufficient. Identify empirical moments that discriminate between legal, informational, and coordination barriers using documented claim structures and settlement histories. This is a restricted mechanism-design target for a stated creditor environment, not a claim that all sovereign restructurings admit one universally optimal procedure.

Give repayment capacity in a resource numeraire and include new debt service after new financing; otherwise borrowing can artificially make every restructuring feasible. Specify a terminal no-Ponzi or repayment condition. Define insolvency relative to the present value of legally feasible repayments under the declared capacity process. For each creditor type \(\theta_j\), write \(U_j^{\cal M}(\theta_j;\hat\theta_j,\hat\theta_{-j})\) for its expected discounted recovery net of enforcement costs. Participation requires \(U_j^{\cal M}(\theta_j;\theta_j,\theta_{-j})\ge U_j^{\rm out}(\theta_j)\); Bayesian truthfulness requires the corresponding inequality in expectation over other types for every report \(\hat\theta_j\). Seniority may be statutory or only a bargaining outcome and must not be silently assumed. Define the crisis-cost objective to include delay, lost output, and financing distortions in common welfare units. This strategic mechanism-design formulation overlaps game theory.
\subsection*{Variants and assumption changes}
The following are \emph{editorial variants}, not additional published open conjectures.
\begin{enumerate}
\item \emph{Complete-information workout.} All claims, capacities, and enforcement options are public. Optimize a repayment and voting procedure subject to legal seniority and creditor participation; the task is constrained settlement design rather than elicitation.
\item \emph{Hidden creditor type and ex ante lending.} Permit private outside options and endogenous initial loan supply. Require truthfulness and participation at settlement and competitive pricing at issuance; compare the ex post relief gain with changed borrowing terms.
\end{enumerate}
\subsection*{References and source audit}
\begin{itemize}
\item Kris James Mitchener; Christoph Trebesch. \emph{Sovereign Debt in the 21st Century}. 2021. Locator: Section8.2, final paragraph, printed p.46; Section9 on holdouts; Section10, printed pp.51--52. Primary text checked. \url{https://www.ifo.de/DocDL/cesifo1_wp8959.pdf}.
\end{itemize}
Official-creditor implications are an agenda; the multilateral mechanism is editorial and overlaps game theory. This formulation contains strategic incentives and therefore overlaps game theory.
\begin{enumerate}\raggedright
\item\hypertarget{definitions:source:30007537:1}{} Kris James Mitchener; Christoph Trebesch. 2021. \emph{Sovereign Debt in the 21st Century}. Section8.2, final paragraph, printed p.46; Section9 on holdouts; Section10, printed pp.51--52.
\url{https://www.ifo.de/DocDL/cesifo1_wp8959.pdf}.
\end{enumerate}

\subsection*{Common conventions: Source mathematical conventions}

These conventions apply to each entry unless explicitly replaced.

\textbf{Models and identification.} A primitive model \(m\) specifies all probability laws, feasible choices, information, equilibrium and measurement rules used by its target. The admissible class \(\mathcal M\) is fixed independently of outcomes. If \(P_O(m)\) is its observable probability law and \(\tau(m)\) the target, its identified set at observable law \(P\) is
\[
 \mathcal I_\tau(P)=\{\tau(m):m\in\mathcal M,\ P_O(m)=P\}.
\]
Point identification means that this set is a singleton. An empty set signals incompatibility of data and assumptions. Coordinate bounds need not describe the joint set. Bounds are sharp only if every retained target is attainable or arbitrarily approachable by compatible models. Finite bounds require explicit restrictions on unobserved quantities; unrestricted confounding, hidden wealth, extrapolation or arbitrary equilibrium selection can yield uninformative sets.

\textbf{Causal interventions.} A potential outcome \(Y_i(p)\) is the outcome under a complete policy regime \(p\), including timing, financing, information and market spillovers. The target population and units are fixed before treatment. With interference, use the complete assignment vector or a justified exposure mapping; individual no-interference notation alone cannot identify an economy-wide intervention. A difference of mean potential outcomes does not require the cross-world joint distribution. Conditioning on post-treatment migration, survival, enrollment or employment changes the estimand and needs separate justification.

\textbf{Statistical statements.} An estimator, confidence set, forecast or test specifies a sampling law, model class and loss/coverage criterion. Uniform \((1-\alpha)\)-coverage means \(\inf_{m\in\mathcal M}\Pr_m\{\tau(m)\in C_n\}\geq1-\alpha\), or an explicitly asymptotic version. The whole parameter space trivially has coverage, so informativeness requires a stated width, risk or power objective. A forecasting improvement is relative to a named benchmark, information set and evaluation population.

\textbf{Equilibrium and optimization.} An equilibrium comprises feasible plans, optimality, beliefs where needed, budget constraints and all market/resource clearing. The primitives or a stated benchmark define each correspondence \(\mathcal E(p;m)\). Nonemptiness is assumed or proved, never inferred just from compact policy sets. A selected-equilibrium target uses a declared selection rule; a robust target quantifies over all permitted equilibria. Use suprema or infima unless attainment is established. An \(\varepsilon\)-optimal policy is feasible and within \(\varepsilon\) of the finite optimum value. Report an empty feasible set separately.

\textbf{Welfare and accounting.} Normative utility, interpersonal normalization, social weights, numeraire, discounting and the use of public receipts are primitives. Behavioral utility need not coincide with the welfare criterion. Household welfare includes ownership income and rents once; adding profits again to compensating variation generally double-counts them. A monetary equivalent is defined by an explicit transfer and price convention, with monotonicity and range conditions. Average effects, mechanism contributions, transfers and welfare losses are different targets. Infinite sums require convergence and dynamic feasible plans require terminal or no-Ponzi conditions.

\textbf{Source versions.} A working-paper date, revision date and journal publication year are distinguished. A DOI for a journal version is not automatically the DOI of an earlier manuscript. Locators refer to the stated version and printed pages where specified. A metadata-only check does not verify a claim about a full-text conclusion. Every entry's source subsection narrows the provenance claim accordingly.
% END ECONOMICS STATEMENT
\end{document}
